% !TeX program = pdflatex
\documentclass[11pt]{article}
\usepackage[T1]{fontenc}
\usepackage{lmodern,booktabs,hyperref}
\usepackage[margin=1in]{geometry}
\title{PasTeX Developer Reference}
\author{PasTeX contributors}
\date{}
\begin{document}
\maketitle

\section{Build}
Install Lazarus with LCL and SynEdit and its matching Free Pascal compiler.
Linux builds support Qt6, Qt5, GTK2 and GTK3. Qt requires the corresponding
QtPas bridge. Select a widget set and build mode as follows.
\begin{verbatim}
./build.sh default qt6
./build.sh debug qt6
./build.sh release qt6
\end{verbatim}
The default widget set is GTK2. Use \texttt{./build.sh release} to build all
four Linux variants. Override \texttt{LAZBUILD} or \texttt{LAZFLAGS} for a
custom Lazarus installation.

On Windows, run the following in PowerShell.
\begin{verbatim}
.\build-windows.ps1 -Mode Default
.\build-windows.ps1 -Mode Release
\end{verbatim}
Open \texttt{pastex.lpi} to work in the Lazarus IDE. Default and debug builds
go to \texttt{bin/default/} and \texttt{bin/debug/}. Release builds use
\texttt{bin/release/vX.X.X/} and include the CPU, OS and widget set in the
executable name.

A \TeX{} distribution enables compiled preview, bibliography processing and
PDF export. Poppler tools or Poppler GLib provide PDF display. FreeType and
HarfBuzz enable outline rendering and OpenType shaping, with widget text as
a fallback. Install the packages and fonts required by the document.

\section{Source structure}
\begin{center}\small
\begin{tabular}{ll}\toprule
Area & Units in \texttt{src/}\\\midrule
Syntax and expansion & paragraphmodel, mathmodel, texmacros\\
Native text and math & draftlayout, texlinebreak, texfontdata\\
Outlines and font selection & nativetexglyphs, texfontshapes, texfontconfig\\
Page layout & nativepagination, nativefloats, pagedview\\
Diagrams & nativediagrams, diagrampreview, plotgeometry\\
References and bibliography & texcontext, nativegenerated, bibliographydata\\
Compiled preview & texrenderer, formularender, vectorpreview\\\bottomrule
\end{tabular}
\end{center}
Native layout uses \TeX{} font metrics and glyph outlines. Paged and continuous
editing share the renderer. The selected \TeX{} engine produces the final PDF.
Implementation references include the
\href{https://tug.org/texlive/doc/tlbuild.pdf}{TeX Live build and source guide},
the \href{https://www.latex-project.org/help/documentation/source2e.pdf}{LaTeX kernel sources},
and the \href{https://ctan.org/pkg/amsmath}{AMS math sources}.

\section{Fonts and artwork}
The executable embeds baseline TeX fonts, metrics and notices in
\texttt{resources/native-fonts.res}. To regenerate the bundle with \TeX{}
installed, run \texttt{python3 scripts/bundle\_native\_fonts.py}. The fonts
retain the licenses listed in \texttt{resources/FONT-LICENSES.txt}.

Open \texttt{src/mainform.pas} in Lazarus and press F12 for the form designer.
Artwork lives in \texttt{icons/}. Regenerate the application and toolbar icon
exports with \texttt{icons/export\_icons.py} and
\texttt{icons/export\_toolbar\_ico.py}.
\end{document}
